% !TEX TS-program = lualatex
Метод временных разниц был разработан в 1980-х годах, когда учёные начали искать более эффективные способы обучения агентов на основе обратной связи от среды. Одной из основных причин его появления стало ограничение метода Монте-Карло, который до этого доминировал в области обучения с подкреплением. Главным недостатком метода Монте-Карло была необходимость завершения всего эпизода (игры), прежде чем можно было бы оценить текущее состояние.

Основная идея классического метода заключается в том, чтобы обновлять оценки состояний, измеряя разницу между текущей и следующей оценкой, при этом учитывая будущие награды. Формально, алгоритм $TD(0)$ --- простейшая форма обучения с разностью времён --- описывается следующим образом:

\begin{equation}\label{td-error}
    \delta_t = R(s_t, a_t) + \gamma V(s_{t+1}) - V(s_t).
\end{equation}

\begin{equation}\label{td-base}
    V(s_t) \leftarrow V(s_t) + \alpha \delta_t.
\end{equation}

\begin{itemize}
    \item $\delta_t$ --- ошибка алгоритма на шаге времени $t$.
    \item $R(s_t, a_t)$ --- награда, полученная после выполнения действия $a_t$ в состоянии $s_t$.
    \item $\gamma$ --- коэффициент дисконтирования, который определяет, насколько важны будущие награды.
    \item $\alpha$ --- скорость обучения, которая контролирует, насколько сильно обновляется текущее значение на основе ошибки алгоритма.
    \item $V(s_t)$ --- оценка состояния $s_t$ (до обновления).
    \item $V(s_{t+1})$ --- оценка следующего состояния $s_{t+1}$.
\end{itemize}

Существенный вклад в развитие этого метода внёс Ричард Саттон, представивший в 1988 году алгоритм $TD(\lambda)$ --- один из первых методов, использующих подход для решения задач обучения с подкреплением.

Алгоритм $TD(\lambda)$ является обобщением уравнения (\ref{td-base}) и включает использование следов посещения. Это расширение позволяет агенту обновлять не только значение текущего состояния, но и значения предыдущих состояний вдоль траектории. Параметр $\lambda$ определяет степень влияния более ранних состояний на процесс обновления.

Правило обновления для $TD(\lambda)$ задается следующим образом:

\begin{equation}\label{td-lambda}
    V(s_t) \leftarrow V(s_t) + \alpha \delta_t \cdot \lambda^{t - k},
\end{equation}
где $\lambda$ --- параметр, который определяет, как далеко в прошлое должны распространяться следы пригодности, $\lambda \in [0, 1]$.

По мере увеличения разности $t - k$ вклад предыдущих состояний уменьшается экспоненциально по формуле $\lambda^{t-k}$. При $t - k = 1$ алгоритм приближается к методу Монте-Карло, а при $t - k = 0$ совпадает со стандартным методом $TD(0)$.
